% !TEX root = ../main.tex
\chapter{Signal Processing Windows}
\label{ch:signal_windows}

Window functions are the unsung companions of the FFT. A raw
finite-length signal, when transformed, exhibits spectral
leakage --- energy from each frequency bin spreads into
neighbouring bins because the implicit periodic extension of the
finite window has discontinuities. Multiplying the signal by a
smooth window (zero at the edges, peak in the middle) eliminates
the discontinuity and confines the leakage to a controlled
side-lobe pattern.

\Lucid\ provides the standard window family at
\code{lucid.signal.windows}, mirroring SciPy's
\code{scipy.signal.windows}. The classic reference for the
trade-offs is Harris's 1978 survey
\cite{harris_blackman_1978}. This chapter describes each window,
the trade-offs between them, and the standard usage patterns.

\section{Why Windowing Matters}
\label{sec:signal_windows:why}

The DFT assumes the input is one period of a periodic signal. A
finite-length real signal is almost never genuinely periodic; the
implied periodic extension has a discontinuity at the wraparound
boundary that the FFT interprets as a broadband signal.

\figref{fig:signal_windows:leakage} sketches the phenomenon for
a pure sinusoid: with no window, the spectrum has substantial
energy in bins other than the sinusoid's frequency (spectral
leakage); with a Hann window, the leakage is confined to a few
neighbouring bins.

\begin{figure}[t]
\centering
\begin{adjustbox}{max width=\textwidth, center}
\begin{tikzpicture}[
  ax/.style={->, >=stealth, thick},
  bin/.style={fill=lucid.info.bg, thick},
  leak/.style={fill=lucid.warn.bg!50, thick},
  hdr/.style={font=\small\sffamily\bfseries, align=center},
]
  % Left: no window
  \node[hdr] at (2, 4) {No window};
  \draw[ax] (0, 0) -- (4.5, 0) node[right, font=\small] {bin};
  \draw[ax] (0, 0) -- (0, 3.5) node[above, font=\small] {$|X|$};
  \foreach \i/\h in {0/0.6, 1/0.9, 2/1.5, 3/3.0, 4/1.5, 5/0.9, 6/0.6, 7/0.4, 8/0.3} {
    \fill[leak] ({0.4 + \i*0.45}, 0) rectangle ({0.7 + \i*0.45}, \h);
  }

  % Right: Hann window
  \node[hdr] at (9, 4) {With Hann window};
  \draw[ax] (6.5, 0) -- (11, 0) node[right, font=\small] {bin};
  \draw[ax] (6.5, 0) -- (6.5, 3.5) node[above, font=\small] {$|X|$};
  \foreach \i/\h in {0/0.05, 1/0.1, 2/0.5, 3/3.0, 4/0.5, 5/0.1, 6/0.05, 7/0.02, 8/0.02} {
    \fill[bin] ({6.9 + \i*0.45}, 0) rectangle ({7.2 + \i*0.45}, \h);
  }
\end{tikzpicture}
\end{adjustbox}
\caption[Spectral leakage with and without windowing.]{Schematic
of spectral leakage. A pure sinusoid sampled into a finite-
length window, then FFT'd: without windowing (left), the
spectrum spreads over many bins; with a Hann window (right), the
spread is confined to a few neighbouring bins. The peak height
is reduced because the window itself reduces total signal
energy, but the spectral concentration is dramatically improved.}
\label{fig:signal_windows:leakage}
\end{figure}

\subsection{The trade-off}
\label{ssec:signal_windows:trade_off}

Windowing trades two things: \emph{main-lobe width} (the spread
of energy from a single frequency component into adjacent bins)
against \emph{side-lobe level} (how far away the spread extends
before fully attenuating).

\begin{itemize}
  \item Rectangular (no window): main lobe is narrowest possible
    (1 bin), side lobes are loudest possible (~-13~dB).
  \item Hann: main lobe is wider (~2 bins), side lobes are
    quieter (~-32~dB).
  \item Blackman: even wider main lobe (~3 bins), even quieter
    side lobes (~-58~dB).
  \item Kaiser-Bessel: parameterised by a scalar that controls
    the trade-off; can approach optimal main-lobe vs side-lobe.
\end{itemize}

The right choice depends on the application: closely spaced
spectral components want narrow main lobes; faint signals next to
loud ones want low side lobes.

\section{Window Functions}
\label{sec:signal_windows:functions}

\Lucid's \code{lucid.signal.windows} subpackage provides eight
standard windows. \tabref{tab:signal_windows:catalog} summarises.

\begin{table}[t]
\centering
\small
\begin{tabular}{lllc}
\toprule
Window & Formula sketch & Main lobe & Side lobe \\
\midrule
rectangular  & $w_n = 1$                       & 1 bin   & -13~dB \\
hann         & $0.5 (1 - \cos(2\pi n/N))$     & 2 bins  & -32~dB \\
hamming      & $0.54 - 0.46 \cos(2\pi n/N)$   & 2 bins  & -43~dB \\
blackman     & sum of 3 cosines               & 3 bins  & -58~dB \\
blackman-harris & sum of 4 cosines            & 3 bins  & -92~dB \\
bartlett     & triangular                     & 2 bins  & -27~dB \\
flattop      & 5-cosine for accurate amplitude & 5 bins & -90~dB \\
kaiser       & Bessel-based, parameterised    & variable & variable \\
\bottomrule
\end{tabular}
\caption[Window function catalogue.]{Standard window functions
exposed by \code{lucid.signal.windows}. The main-lobe width and
side-lobe level are the standard characterisations. Kaiser is
parameterised by $\beta$, which interpolates between
rectangular ($\beta = 0$) and very wide low-leakage windows
($\beta \to \infty$).}
\label{tab:signal_windows:catalog}
\end{table}

\subsection{Hann}
\label{ssec:signal_windows:hann}

\code{hann(M, sym=True)} returns a length-$M$ Hann window:
\[
w_n \;=\; \tfrac{1}{2}\left(1 - \cos \tfrac{2\pi n}{M - 1}\right),
\qquad n = 0, 1, \ldots, M-1.
\]
Equivalently $w_n = \sin^2(\pi n / (M-1))$. The window is zero at
both endpoints, smoothly rising to a peak of $1$ at the centre.

The Hann is the default for STFT-based spectrogram computation:
the time-frequency trade-off is balanced for typical audio
applications.

\subsection{Hamming}
\label{ssec:signal_windows:hamming}

Similar to Hann but with slightly different coefficients:
\[
w_n \;=\; 0.54 - 0.46 \cos \tfrac{2\pi n}{M-1}.
\]
The endpoints are $w_0 = w_{M-1} = 0.08$ rather than $0$, which
slightly raises the time-domain DC offset but reduces the first
side-lobe by about $10$~dB compared to Hann. The trade-off is
application-dependent.

\subsection{Blackman family}
\label{ssec:signal_windows:blackman}

\code{blackman(M)} is the three-term cosine sum:
\[
w_n \;=\; 0.42 - 0.5 \cos \tfrac{2\pi n}{M-1}
       + 0.08 \cos \tfrac{4\pi n}{M-1}.
\]
Lower side lobes than Hann (about $-58$~dB vs $-32$~dB) at the
cost of a wider main lobe.

\code{blackman\_harris(M)} is a four-term sum:
\[
w_n \;=\; a_0 - a_1 \cos \tfrac{2\pi n}{M-1}
       + a_2 \cos \tfrac{4\pi n}{M-1}
       - a_3 \cos \tfrac{6\pi n}{M-1},
\]
with $(a_0, a_1, a_2, a_3) \approx (0.3585, 0.4884, 0.1413,
0.0118)$. Side lobes drop to $-92$~dB, useful for detecting very
faint signals next to loud ones.

\subsection{Bartlett}
\label{ssec:signal_windows:bartlett}

The triangular window: linear ramp up, linear ramp down. Simple
to compute, modest performance.

\subsection{Kaiser}
\label{ssec:signal_windows:kaiser}

\code{kaiser(M, beta)}: parameterised by $\beta$, the Bessel-
function shape parameter:
\[
w_n = \frac{I_0\left(\beta \sqrt{1 - \left(\frac{2n}{M-1} -
1\right)^2}\right)}{I_0(\beta)},
\]
where $I_0$ is the zeroth-order modified Bessel function.

Larger $\beta$ gives wider main lobe and lower side lobes. The
common choices: $\beta = 5$ (similar to Hamming), $\beta = 8.6$
(similar to Blackman-Harris).

\section{Symmetric vs Periodic Windows}
\label{sec:signal_windows:sym_periodic}

A subtle but important convention: window functions can be
generated in \emph{symmetric} form (designed for filtering, where
symmetric means linear phase) or \emph{periodic} form (designed
for STFT-like overlap-add, where periodic means the wraparound
is seamless).

The two differ only at the endpoints: the symmetric form has
$w_0 = w_{M-1}$ (or both zero); the periodic form has $w_0 \neq
w_{M-1}$ (specifically, the periodic form is one period of an
extension that would wrap around).

\Lucid's windows take a \code{sym} parameter (default
\code{True}). For STFT use, set \code{sym=False}.

\section{Worked Example: A Mel-Spectrogram Frontend}
\label{sec:signal_windows:mel_example}

A complete audio frontend that converts a raw waveform to a
log-mel spectrogram:

\begin{lstlisting}[style=lucid,language=LucidPython]
import lucid
import lucid.fft as F
import lucid.signal.windows as W

def log_mel_spectrogram(audio, sr=16000, n_fft=1024,
                        hop_size=256, n_mels=80):
    # Frame the signal
    window = W.hann(n_fft, sym=False)
    n_frames = (audio.shape[0] - n_fft) // hop_size + 1
    frames = lucid.stack([
        audio[i*hop_size:i*hop_size + n_fft] * window
        for i in range(n_frames)
    ])
    # STFT magnitude
    spec = F.rfft(frames, dim=-1).abs()
    # Mel filterbank (precomputed)
    mel_filters = _build_mel_filterbank(n_fft, sr, n_mels)
    mel_spec = spec @ mel_filters.T
    # Log scale
    return (mel_spec + 1e-10).log()
\end{lstlisting}

The pipeline is the standard for speech recognition front-ends:
window the audio, FFT, project onto mel-scaled filters, take the
log. Every step uses \Lucid\ primitives; the gradient flows back
through the whole chain.

\section{Implementation}
\label{sec:signal_windows:implementation}

The windows live in \code{lucid/signal/windows.py}. Each is a
short composite (\chref{ch:composite_ops}) built from arange,
cosine, and arithmetic. The windows are numpy-free.

The windows are not differentiable in any useful sense (they
are deterministic functions of the size $M$, not of input
tensors), so they have no backward.

\subsection{Caching}
\label{ssec:signal_windows:caching}

A typical training loop calls the same window with the same size
many times per step. \Lucid\ does not internally cache the
window, but the user can cache it once and reuse across calls.
The CPU cost of generating a window is microseconds and rarely
worth caching for performance; the cache is for code clarity.

\subsection{Length conventions}
\label{ssec:signal_windows:length}

The window length $M$ is the FFT length $N$. The window is
multiplied elementwise with the signal frame; if the signal is
zero-padded (common for non-power-of-2 frame sizes), the window
should be applied before padding.

\section{Window-Function Backward Through STFT}
\label{sec:signal_windows:stft_backward}

Although the windows themselves are not differentiable, the
STFT pipeline that uses them is. The window appears as a
constant in the chain (\code{frames * window}); the backward
treats it as a constant multiplier and propagates the upstream
cotangent multiplied by the window.

The user-side consequence: gradient flows naturally through
window-based pipelines. A loss computed on a spectrogram
backwards to the raw audio without special treatment.

\section{Choosing a Window}
\label{sec:signal_windows:choosing}

A short guide.

\begin{itemize}
  \item \textbf{Default audio STFT}: Hann.
  \item \textbf{Closely-spaced spectral components}: rectangular
    (narrow main lobe) or Hamming.
  \item \textbf{Faint signals next to loud}: Blackman, Blackman-
    Harris, or Kaiser with high $\beta$.
  \item \textbf{Accurate amplitude measurement of a known
    spectral component}: flat-top.
  \item \textbf{Linear-phase filter design}: any symmetric form;
    Hamming or Blackman.
\end{itemize}

For most ML applications (speech recognition, audio
classification, music generation), Hann is the default and works
well.

\section{Spectral Leakage: A Quantitative Analysis}
\label{sec:signal_windows:leakage_quant}

The leakage pattern of a window function is its Fourier transform
$W(\omega)$. For a sinusoid of frequency $\omega_0$ multiplied by
the window $w$ and transformed, the output spectrum is the
convolution
\[
\widehat{(w \cdot x)}(\omega)
\;=\; \int W(\omega - \omega') \cdot \delta(\omega' - \omega_0)\, d\omega'
\;=\; W(\omega - \omega_0).
\]
The window's Fourier transform $W$ \emph{is} the spectral leakage
pattern, shifted to the sinusoid's frequency.

\subsection{The rectangular case}
\label{ssec:signal_windows:rect_dft}

The rectangular window of length $N$ has Fourier transform
\[
W_{\mathrm{rect}}(\omega) \;=\; \frac{\sin(N\omega/2)}{\sin(\omega/2)},
\]
the \emph{Dirichlet kernel}. The main lobe has width $4\pi/N$
(zero crossings at $\omega = \pm 2\pi/N$); the side lobes have
peak amplitude $1/(\sin(\pi \cdot 1.5/N) \cdot N) \approx
2/(3\pi)$ relative to the main lobe, giving the famous
$-13.3$\,dB attenuation.

\subsection{The Hann case}
\label{ssec:signal_windows:hann_dft}

The Hann window is a cosine raised to the first power:
\[
w_n \;=\; \tfrac{1}{2}(1 - \cos 2\pi n / N).
\]
Its Fourier transform is a sum of three Dirichlet kernels
(at $\omega = 0, \pm 2\pi/N$), each weighted by the cosine
coefficients. The result has a main lobe twice as wide as the
rectangular window but side lobes attenuated by an additional
$\sim 20$\,dB.

The trade-off generalises: an $M$-cosine window has $M$ times
the rectangular main-lobe width and side lobes attenuated by
roughly $20 (M-1)$\,dB.

\section{The COLA Property}
\label{sec:signal_windows:cola}

For STFT to be invertible (the inverse short-time Fourier
transform reconstructs the original signal exactly), the window
must satisfy the \emph{constant overlap-add} (COLA) property:
\[
\sum_{m=-\infty}^{\infty} w(n - mH) \;=\; \text{constant},
\]
where $H$ is the hop size. The condition says that the summed
overlap of shifted windows is uniform across the time axis.

\subsection{COLA-compliant configurations}
\label{ssec:signal_windows:cola_configs}

Common COLA-compliant choices:
\begin{itemize}
  \item Hann window with $H = N/2$: $\sum_m w(n - mH) = 1$.
  \item Hann window with $H = N/4$: $\sum = 1.5$.
  \item Hamming window with $H = N/2$: $\sum = 1.08$.
  \item Rectangular window with $H = N$: trivially COLA (no
    overlap).
\end{itemize}

\subsection{Reconstructing from STFT}
\label{ssec:signal_windows:istft}

Given a STFT $X[m, k]$ produced with window $w$ and hop $H$, the
inverse STFT reconstructs
\[
\hat x[n] \;=\; \frac{\sum_m w(n - mH) \cdot
\mathrm{ifft}(X[m, :])[n - mH]}{\sum_m w(n - mH)^2}.
\]
The denominator is the COLA-squared term. When the window has
been chosen COLA-compliant and the data unaltered, the
reconstruction is exact (modulo edge effects at the very start
and end of the signal).

\section{Window Optimisation: The Kaiser--Bessel Family}
\label{sec:signal_windows:kaiser_opt}

The Kaiser window \cite{kaiser_window_1980} provides the closest
approximation to the \emph{prolate spheroidal wave function},
which is optimal in the sense of maximising the fraction of
energy in the main lobe for a given side-lobe level.

\subsection{The Kaiser--Bessel formula}
\label{ssec:signal_windows:kaiser_formula}

\[
w_n \;=\; \frac{I_0\!\left(\beta \sqrt{1 - \left(\dfrac{2n}{N-1} - 1\right)^2}\right)}{I_0(\beta)},
\quad n = 0, 1, \ldots, N-1,
\]
where $I_0$ is the zeroth-order modified Bessel function
(\chref{ch:special_functions}). The parameter $\beta$ controls
the trade-off:

\begin{itemize}
  \item $\beta = 0$: $I_0(0) = 1$ everywhere; the window is
    rectangular.
  \item $\beta = 5.0$: side lobes around $-37$\,dB; similar to
    Hamming.
  \item $\beta = 8.6$: side lobes around $-65$\,dB; similar to
    Blackman.
  \item $\beta = 12.0$: side lobes around $-100$\,dB.
\end{itemize}

\subsection{Choosing beta}
\label{ssec:signal_windows:kaiser_beta}

For a desired side-lobe attenuation $A$ in dB:
\[
\beta \;\approx\;
\begin{cases}
0 & A \leq 21, \\
0.5842 (A - 21)^{0.4} + 0.07886 (A - 21) & 21 < A \leq 50, \\
0.1102 (A - 8.7) & A > 50.
\end{cases}
\]
The corresponding main-lobe width is
\[
\Delta \omega \;\approx\; \frac{2\pi (A - 8)}{2.285 (N - 1)}.
\]
A wider main lobe trades for lower side lobes; the user picks
the balance.

\section{Window Energy Corrections}
\label{sec:signal_windows:energy_corrections}

When a window is applied before FFT, the resulting spectrum is
underestimated relative to the un-windowed signal because the
window reduces total energy. Two correction conventions apply.

\subsection{Coherent gain}
\label{ssec:signal_windows:coherent_gain}

The mean of the window:
\[
G_{\mathrm{coh}} \;=\; \frac{1}{N} \sum_n w_n.
\]
Used to correct the amplitude of a known spectral component.
Hann: $G_{\mathrm{coh}} = 0.5$. Hamming: $G_{\mathrm{coh}} = 0.54$.

\subsection{Power gain}
\label{ssec:signal_windows:power_gain}

The mean of the squared window:
\[
G_{\mathrm{pow}} \;=\; \frac{1}{N} \sum_n w_n^2.
\]
Used to correct the power of a broadband signal. Hann:
$G_{\mathrm{pow}} = 3/8$. Hamming: $G_{\mathrm{pow}} \approx
0.397$.

These corrections appear inside the Welch PSD estimator
(\chref{ch:fft}); the user does not see them directly in
\Lucid's API but must apply them if interpreting raw STFT
magnitudes as physical units.

\section{Applications Beyond STFT}
\label{sec:signal_windows:applications}

Window functions appear outside spectrogram computation.

\subsection{Pre-emphasis filtering}
\label{ssec:signal_windows:pre_emphasis}

Speech recognition often applies a first-order high-pass filter
$y_n = x_n - 0.97 x_{n-1}$ before windowing. The
combined filter is a window-equivalent; \Lucid\ does not provide
a separate primitive but the user can chain the filter and the
window manually.

\subsection{Filter design}
\label{ssec:signal_windows:filter_design}

FIR filter design via the window method: take the inverse
Fourier transform of the desired frequency response, truncate to
a finite length, and apply a window to reduce ripple. The chosen
window controls the trade-off between transition bandwidth and
stop-band attenuation.

\subsection{Image processing}
\label{ssec:signal_windows:image_proc}

2-D windows for image processing apply the same principles
separately along each axis. \Lucid\ does not provide 2-D windows
directly; the user can compute an outer product:
\code{w2 = w.unsqueeze(0) * w.unsqueeze(1)}.

\section{Summary and Forward References}
\label{sec:signal_windows:summary}

\code{lucid.signal.windows} provides the standard family of
window functions for STFT and spectral analysis. The choice of
window controls the main-lobe / side-lobe trade-off; Hann is the
typical default. Windows are not differentiable themselves but
participate transparently in differentiable STFT pipelines.

The next chapter, \chref{ch:special_functions}, treats the
special-function subpackage (erf, erfinv, Bessel, etc.).
\chref{ch:distributions} treats probability distributions.
\chref{ch:einops} closes \partref{part:math_subsystems} with
tensor rearrangement algebra.

\begin{lucidnote}
For the user: \code{lucid.signal.windows.hann(N)} is the
typical entry point. The result is a 1-D tensor of the requested
length, ready to multiply with each frame of an STFT.
\end{lucidnote}
